\documentclass[12pt]{article}
\usepackage{amsmath, amssymb, amsthm}
\usepackage{geometry}
\geometry{letterpaper, margin=1in}
\usepackage{hyperref}
\usepackage{booktabs}
\usepackage{enumitem}

\title{Math Problem Sets: Characterizations of Invertible Matrices}
\author{}
\date{}

\begin{document}

\maketitle

\tableofcontents
\newpage

\section{Concept Checklist \& Question Mapping}
\label{sec:checklist_start}
This section outlines the core concepts and problem types extracted from the provided text, ensuring complete coverage across the problem sets.

\begin{itemize}
    \item \textbf{C1: $2 \times 2$ Matrix Invertibility.} Calculating determinants ($ad-bc$) to check invertibility and using the 2D inverse formula. (\textit{Questions: 1, 2, 19, 20})
    \item \textbf{C2: Triangular Matrix Tricks.} Identifying upper/lower triangular matrices and checking their main diagonals for zeros. (\textit{Questions: 3, 8, 17})
    \item \textbf{C3: Recognizing Linear Dependence.} Spotting identical columns/rows, zero columns, or obvious scalar multiples without full row reduction. (\textit{Questions: 4, 5, 6, 7})
    \item \textbf{C4: IMT - Null Space \& Solutions.} Connecting trivial null spaces ($Ax=0$) and inconsistencies to invertibility via the Invertible Matrix Theorem (IMT). (\textit{Questions: 9, 11, 14, 22})
    \item \textbf{C5: IMT - Span \& Linear Independence.} Relating spanning $\mathbb{R}^n$ and independent columns to invertibility. (\textit{Questions: 10, 13, 16})
    \item \textbf{C6: IMT - Linear Transformations.} Equating one-to-one and onto mappings for square matrices. (\textit{Questions: 21, 29})
    \item \textbf{C7: Matrix Operations \& Properties.} Applying product inverse, transpose inverse, and associativity rules ($(AB)^{-1} = B^{-1}A^{-1}$). (\textit{Questions: 15, 18, 23, 28})
    \item \textbf{C8: Theoretical Proofs.} Leveraging algebraic manipulations and IMT equivalences to prove invertibility of complex matrix expressions. (\textit{Questions: 24, 25, 30})
    \item \textbf{C9: Block Matrix Inversion.} Applying the partitioned block matrix formula for inverses. (\textit{Questions: 26, 27})
    \item \textbf{C10: Floating-Point / Condition Number.} Understanding the practical computational limits of mathematically invertible matrices. (\textit{Questions: 31})
\end{itemize}

\newpage

\section{Problem Sets}
\label{sec:problems}

\subsection*{Part A: Invertibility Checks}
\textit{Objective: Determine if the following matrices are invertible with minimal calculation. Provide a brief justification.}

\begin{enumerate}[label=\textbf{\arabic*.}]
    \item $A = \begin{bmatrix} 3 & 5 \\ 2 & 4 \end{bmatrix}$
    \item $B = \begin{bmatrix} -6 & 9 \\ 2 & -3 \end{bmatrix}$
    \item $C = \begin{bmatrix} 4 & 0 & 0 \\ 7 & -2 & 0 \\ 1 & 5 & 3 \end{bmatrix}$
    \item $D = \begin{bmatrix} 8 & -1 & 4 \\ 0 & 0 & 0 \\ 2 & 9 & 3 \end{bmatrix}$
    \item $E = \begin{bmatrix} 1 & -3 & 2 \\ -2 & 6 & -4 \\ 5 & 1 & 8 \end{bmatrix}$
    \item $F = \begin{bmatrix} 1 & 4 & 5 \\ 2 & 1 & 3 \\ 3 & 5 & 8 \end{bmatrix}$
    \item $G = \begin{bmatrix} 2 & 0 & 3 \\ 0 & 1 & 4 \\ 0 & 0 & 0 \end{bmatrix}$
    \item $H = \begin{bmatrix} -1 & 5 & 9 & 2 \\ 0 & 3 & 4 & 8 \\ 0 & 0 & -2 & 7 \\ 0 & 0 & 0 & 4 \end{bmatrix}$
\end{enumerate}

\subsection*{Part B: Implications of the Invertible Matrix Theorem}
\textit{Objective: Determine if the statements are True or False, or answer the question. Justify based on the IMT.}

\begin{enumerate}[label=\textbf{\arabic*.}, resume]
    \item \textbf{(T/F)} If an $n \times n$ matrix $A$ has a trivial null space ($Ax = 0$ has only the trivial solution), then the columns of $A$ span $\mathbb{R}^n$.
    \item \textbf{(T/F)} If an $n \times n$ matrix $B$ cannot be row reduced to the identity matrix $I_n$, then the equation $Bx = b$ has a solution for all $b \in \mathbb{R}^n$.
    \item \textbf{(T/F)} If $C$ is an $n \times n$ matrix and $Cx = 0$ has a non-trivial solution, then the linear transformation $x \mapsto Cx$ is one-to-one.
    \item \textbf{(T/F)} A $4 \times 5$ matrix whose columns span $\mathbb{R}^4$ is invertible.
    \item \textbf{(T/F)} If the columns of a square matrix $D$ are linearly independent, then $D^T$ is invertible.
    \item If $E$ is an $n \times n$ matrix and $Ex = b$ is inconsistent for some $b \in \mathbb{R}^n$, what can be said about the number of pivots in $E$?
    \item Let $F$ be an $n \times n$ matrix. If there exists a matrix $G$ such that $FG = I_n$, must $GF = I_n$?
    \item If $H$ is a $6 \times 6$ matrix and the equation $Hx = c$ has unique solutions for every $c$, do the columns of $H$ form a basis for $\mathbb{R}^6$?
    \item \textbf{(T/F)} When a square upper triangular matrix has non-zero entries on its main diagonal, its columns are linearly independent.
    \item If matrix $M$ is invertible, is $M^3$ guaranteed to be invertible?
\end{enumerate}

\subsection*{Part C: Linear Transformations}
\textit{Objective: Extract standard matrices and utilize IMT logic for transformations.}

\begin{enumerate}[label=\textbf{\arabic*.}, resume]
    \item Find the standard matrix $A$ for the transformation $T(x_1, x_2) = (4x_1 - x_2, 7x_1 - 2x_2)$. Determine if it is invertible, and if so, find $T^{-1}(x_1, x_2)$.
    \item Show whether the linear transformation $T(x_1, x_2) = (-6x_1 + 8x_2, 3x_1 - 4x_2)$ is invertible.
    \item Suppose a linear transformation $T: \mathbb{R}^n \to \mathbb{R}^n$ is onto. Does the equation $T(x) = 0$ have non-trivial solutions?
\end{enumerate}

\subsection*{Part D: Theoretical Proofs}
\textit{Objective: Leverage associativity and matrix definitions to write formal proofs.}

\begin{enumerate}[label=\textbf{\arabic*.}, resume]
    \item Show that if $A$ is an invertible square matrix and $AB = 0$, then $B = 0$.
    \item If $P$ is an orthogonal matrix (meaning $P^T P = I_n$), prove that $P$ is invertible and that its inverse is $P^T$.
    \item Prove that if $A^2 - 2A + I_n = 0$, then $A$ is invertible.
    \item Prove: If $A, B$, and $C$ are $n \times n$ matrices such that $ABC = I_n$, show that $B$ is invertible.
\end{enumerate}

\subsection*{Part E: Common Variations (Block Matrices)}
\textit{Objective: Apply block matrix inversion rules.}

\begin{enumerate}[label=\textbf{\arabic*.}, resume]
    \item Given $M = \begin{bmatrix} I_3 & C \\ 0 & I_3 \end{bmatrix}$, find $M^{-1}$.
    \item Find the inverse of the partitioned matrix $M = \begin{bmatrix} A & B \\ 0 & D \end{bmatrix}$ where $A = \begin{bmatrix} 3 & 1 \\ 2 & 1 \end{bmatrix}$, $B = \begin{bmatrix} 1 & 0 \\ 0 & 1 \end{bmatrix}$, and $D = \begin{bmatrix} 2 & 0 \\ 0 & 4 \end{bmatrix}$.
\end{enumerate}

\subsection*{Part F: Advanced Diagnostics ("Find the Flaw")}
\textit{Objective: Act as a grader. Locate the subtle mathematical error, explain why it is wrong, and correct it.}

\begin{enumerate}[label=\textbf{\arabic*.}, resume]
    \item \textbf{Problem:} Determine if $A = \begin{bmatrix} 5 & 4 \\ 2 & 3 \end{bmatrix}$ is invertible and find the inverse. \\
    \textbf{Flawed Solution:} The determinant is $\det(A) = (5)(3) + (4)(2) = 15 + 8 = 23$. Since it's not zero, $A$ is invertible. The inverse is $A^{-1} = \frac{1}{23} \begin{bmatrix} 3 & -4 \\ -2 & 5 \end{bmatrix}$.
    
    \item \textbf{Problem:} Simplify the expression $(A^T B^T)^{-1}$ for invertible $n \times n$ matrices. \\
    \textbf{Flawed Solution:} Distribute the inverse: $(A^T B^T)^{-1} = (A^T)^{-1} (B^T)^{-1}$. Since transpose and inverse commute, this simplifies to $(A^{-1})^T (B^{-1})^T$.
    
    \item \textbf{Problem:} Suppose $A$ is a $5 \times 4$ matrix with linearly independent columns. Does $A$ span $\mathbb{R}^5$? \\
    \textbf{Flawed Solution:} Because the columns are linearly independent, the equation $Ax=0$ has only the trivial solution. By the IMT, a matrix with only the trivial solution spans its target space. Therefore, yes, it spans $\mathbb{R}^5$.
    
    \item \textbf{Problem:} Prove that if $A^3 = 0$ for a square matrix $A$, then $A$ is not invertible. \\
    \textbf{Flawed Solution:} Assume $A^3 = 0$. Multiply both sides by $A^{-3}$ to get $A^{-3}A^3 = A^{-3}0$. This simplifies to $I = 0$. Since the identity matrix cannot equal the zero matrix, $A$ cannot be invertible.
\end{enumerate}

\subsection*{Part G: Numerical and Condition Number}
\begin{enumerate}[label=\textbf{\arabic*.}, resume]
    \item If a $6 \times 6$ covariance matrix $\Sigma$ calculated in Python has a condition number $\kappa = 10^{16}$, it is technically invertible mathematically (determinant is not exactly zero). Explain why an asset manager should NOT rely on $\Sigma^{-1}$ to calculate portfolio weights.
\end{enumerate}

\newpage

\section{Detailed Solutions}
\label{sec:solutions}

\textbf{1.} $A = \begin{bmatrix} 3 & 5 \\ 2 & 4 \end{bmatrix}$. $\det(A) = (3)(4) - (5)(2) = 12 - 10 = 2$. Since $2 \neq 0$, it is \textbf{invertible}.

\textbf{2.} $B = \begin{bmatrix} -6 & 9 \\ 2 & -3 \end{bmatrix}$. $\det(B) = (-6)(-3) - (9)(2) = 18 - 18 = 0$. It is \textbf{not invertible}.

\textbf{3.} $C$ is a lower triangular matrix. Its determinant is the product of its diagonal entries: $(4)(-2)(3) = -24 \neq 0$. It is \textbf{invertible}.

\textbf{4.} $D$ contains an entire row of zeros (Row 2). The determinant is 0. It is \textbf{not invertible}.

\textbf{5.} In matrix $E$, Row 2 is exactly $-2$ times Row 1 ($-2 \times [1, -3, 2] = [-2, 6, -4]$). The rows are linearly dependent. \textbf{Not invertible}.

\textbf{6.} In matrix $F$, Row 3 is the sum of Row 1 and Row 2. Because of this linear dependence, it will not have 3 pivots. \textbf{Not invertible}.

\textbf{7.} $G$ is an upper triangular matrix with a 0 on the main diagonal. Its determinant is $2 \times 1 \times 0 = 0$. \textbf{Not invertible}.

\textbf{8.} $H$ is a $4 \times 4$ upper triangular matrix. The determinant is the product of the main diagonal: $(-1)(3)(-2)(4) = 24 \neq 0$. \textbf{Invertible}.

\textbf{9. True.} By the IMT, a trivial null space implies the matrix is invertible, which in turn implies its columns span $\mathbb{R}^n$.

\textbf{10. False.} If it cannot row reduce to $I_n$, it is not invertible. By the IMT, an un-invertible matrix does NOT span $\mathbb{R}^n$, meaning there are some vectors $b$ for which $Bx=b$ has no solution.

\textbf{11. False.} A non-trivial solution to $Cx = 0$ means free variables exist, so information is lost. The mapping cannot be one-to-one.

\textbf{12. False.} The IMT strictly applies only to square ($n \times n$) matrices. A $4 \times 5$ matrix cannot be invertible because it maps between spaces of different dimensions.

\textbf{13. True.} Independent columns mean $D$ is invertible (by IMT). The IMT explicitly states $D$ is invertible if and only if $D^T$ is invertible.

\textbf{14.} If $Ex = b$ is inconsistent for some $b$, $E$ does not span $\mathbb{R}^n$. By the IMT, it is not invertible, which means it has \textbf{fewer than $n$ pivot positions}.

\textbf{15. True.} By the IMT, if a square matrix has a right inverse ($FG = I_n$), then that matrix is its true inverse, meaning it also works as a left inverse ($GF = I_n$).

\textbf{16. Yes, they form a basis.} Unique solutions for every $c$ means $H$ is invertible. By IMT, invertible matrices have linearly independent columns that span $\mathbb{R}^n$, which is the exact definition of a basis.

\textbf{17. True.} Non-zero diagonal entries on a triangular matrix guarantee $n$ pivots, meaning it is invertible. By IMT, an invertible matrix has linearly independent columns.

\textbf{18. True.} The product of invertible matrices is invertible. $M^3 = M \cdot M \cdot M$. Since $M$ is invertible, the product is invertible.

\textbf{19.} Standard matrix $A = \begin{bmatrix} 4 & -1 \\ 7 & -2 \end{bmatrix}$. $\det(A) = -8 - (-7) = -1$. $A^{-1} = \frac{1}{-1} \begin{bmatrix} -2 & 1 \\ -7 & 4 \end{bmatrix} = \begin{bmatrix} 2 & -1 \\ 7 & -4 \end{bmatrix}$. \\
Thus, $T^{-1}(x_1, x_2) = (2x_1 - x_2, 7x_1 - 4x_2)$.

\textbf{20.} Standard matrix $A = \begin{bmatrix} -6 & 8 \\ 3 & -4 \end{bmatrix}$. $\det(A) = 24 - 24 = 0$. Since $\det(A) = 0$, the transformation is \textbf{not invertible}.

\textbf{21.} No. If $T$ is onto, its standard matrix is invertible (by IMT). If the matrix is invertible, it is also one-to-one, meaning $T(x) = 0$ has \textbf{only the trivial solution}.

\textbf{22.} Assume $A$ is invertible, so $A^{-1}$ exists. Multiply both sides of $AB = 0$ on the left by $A^{-1}$: \\
$A^{-1}(AB) = A^{-1}0 \implies (A^{-1}A)B = 0 \implies IB = 0 \implies B = 0$.

\textbf{23.} The definition of an inverse is a matrix $P^{-1}$ such that $P^{-1}P = I_n$. We are given $P^T P = I_n$. By the IMT, a left inverse for a square matrix acts as the true inverse. Thus, $P$ is invertible and $P^{-1} = P^T$.

\textbf{24.} We can factor the equation to isolate $I_n$. \\
$I_n = -A^2 + 2A \implies I_n = A(-A + 2I_n)$. \\
Since there exists a matrix $B = (-A + 2I_n)$ such that $AB = I_n$, $A$ is invertible by the IMT.

\textbf{25.} By associativity, $(AB)C = I_n \implies A(BC) = I_n$. By the IMT, since $A$ multiplied by $(BC)$ yields $I_n$, $A$ is invertible. Let's rewrite as $(AB)C = I_n$. This shows $C$ is invertible and $C^{-1} = AB$. Similarly, $A(BC) = I_n \implies A^{-1} = BC$. Since $A$ and $C$ are invertible, and $B = A^{-1} C^{-1}$, $B$ is the product of two invertible matrices, hence $B$ is invertible. 

\textbf{26.} For $M = \begin{bmatrix} I_3 & C \\ 0 & I_3 \end{bmatrix}$, the inverse formula is $M^{-1} = \begin{bmatrix} I_3^{-1} & -I_3^{-1} C I_3^{-1} \\ 0 & I_3^{-1} \end{bmatrix} = \begin{bmatrix} I_3 & -C \\ 0 & I_3 \end{bmatrix}$.

\textbf{27.} First, invert the diagonal blocks. \\
$A^{-1} = \frac{1}{3-2} \begin{bmatrix} 1 & -1 \\ -2 & 3 \end{bmatrix} = \begin{bmatrix} 1 & -1 \\ -2 & 3 \end{bmatrix}$. \\
$D^{-1} = \begin{bmatrix} 1/2 & 0 \\ 0 & 1/4 \end{bmatrix}$. \\
Calculate the top-right block: $-A^{-1} B D^{-1}$. Since $B = I_2$, this is just $-A^{-1} D^{-1}$. \\
$-\begin{bmatrix} 1 & -1 \\ -2 & 3 \end{bmatrix} \begin{bmatrix} 1/2 & 0 \\ 0 & 1/4 \end{bmatrix} = -\begin{bmatrix} 1/2 & -1/4 \\ -1 & 3/4 \end{bmatrix} = \begin{bmatrix} -1/2 & 1/4 \\ 1 & -3/4 \end{bmatrix}$. \\
Assembling the block matrix yields $M^{-1} = \begin{bmatrix} 1 & -1 & -1/2 & 1/4 \\ -2 & 3 & 1 & -3/4 \\ 0 & 0 & 1/2 & 0 \\ 0 & 0 & 0 & 1/4 \end{bmatrix}$.

\textbf{28.} \textbf{Flaw:} The determinant formula is $ad - bc$, not $ad + bc$. \\
\textbf{Correction:} $\det(A) = (5)(3) - (4)(2) = 15 - 8 = 7$. The correct inverse is $A^{-1} = \frac{1}{7} \begin{bmatrix} 3 & -4 \\ -2 & 5 \end{bmatrix}$.

\textbf{29.} \textbf{Flaw:} The inverse of a matrix product requires reversing the order of the terms (the "Socks and Shoes" rule). $(XY)^{-1} \neq X^{-1}Y^{-1}$. \\
\textbf{Correction:} $(A^T B^T)^{-1} = (B^T)^{-1} (A^T)^{-1} = (B^{-1})^T (A^{-1})^T$.

\textbf{30.} \textbf{Flaw:} The student applied the Invertible Matrix Theorem to a non-square matrix. The IMT ONLY applies to $n \times n$ matrices. \\
\textbf{Correction:} A $5 \times 4$ matrix physically cannot span $\mathbb{R}^5$ because it only has 4 column vectors, and a minimum of 5 linearly independent vectors is required to span a 5D space.

\textbf{31.} \textbf{Flaw:} The proof uses circular logic. By multiplying by $A^{-3}$ (the inverse of $A^3$), the student inherently assumed $A$ was invertible in order to prove it was not invertible. \\
\textbf{Correction:} Use determinants. If $A^3 = 0$, then $\det(A^3) = \det(0) = 0$. Since $\det(A^3) = (\det(A))^3$, we have $(\det(A))^3 = 0 \implies \det(A) = 0$. Since the determinant is zero, $A$ is not invertible.

\textbf{32 (Number 31).} A condition number of $10^{16}$ signifies an extremely ill-conditioned matrix (its columns are nearly parallel). In finite floating-point computing environments (like Python), infinitesimal round-off errors will explode during the inversion process. The resulting weights calculated via $\Sigma^{-1}$ will be pure numerical noise, destroying the portfolio.

\newpage

\section{Concept Checklist \& Question Mapping}
\label{sec:checklist_end}
\textit{This checklist is identical to the one at the beginning of the document to allow for easy back-referencing of missed concepts.}

\begin{itemize}
    \item \textbf{C1: $2 \times 2$ Matrix Invertibility.} Calculating determinants ($ad-bc$) to check invertibility and using the 2D inverse formula. (\textit{Questions: 1, 2, 19, 20})
    \item \textbf{C2: Triangular Matrix Tricks.} Identifying upper/lower triangular matrices and checking their main diagonals for zeros. (\textit{Questions: 3, 8, 17})
    \item \textbf{C3: Recognizing Linear Dependence.} Spotting identical columns/rows, zero columns, or obvious scalar multiples without full row reduction. (\textit{Questions: 4, 5, 6, 7})
    \item \textbf{C4: IMT - Null Space \& Solutions.} Connecting trivial null spaces ($Ax=0$) and inconsistencies to invertibility via the Invertible Matrix Theorem (IMT). (\textit{Questions: 9, 11, 14, 22})
    \item \textbf{C5: IMT - Span \& Linear Independence.} Relating spanning $\mathbb{R}^n$ and independent columns to invertibility. (\textit{Questions: 10, 13, 16})
    \item \textbf{C6: IMT - Linear Transformations.} Equating one-to-one and onto mappings for square matrices. (\textit{Questions: 21, 29})
    \item \textbf{C7: Matrix Operations \& Properties.} Applying product inverse, transpose inverse, and associativity rules ($(AB)^{-1} = B^{-1}A^{-1}$). (\textit{Questions: 15, 18, 23, 28})
    \item \textbf{C8: Theoretical Proofs.} Leveraging algebraic manipulations and IMT equivalences to prove invertibility of complex matrix expressions. (\textit{Questions: 24, 25, 30})
    \item \textbf{C9: Block Matrix Inversion.} Applying the partitioned block matrix formula for inverses. (\textit{Questions: 26, 27})
    \item \textbf{C10: Floating-Point / Condition Number.} Understanding the practical computational limits of mathematically invertible matrices. (\textit{Questions: 31})
\end{itemize}

\end{document}